% Per-game replacements of the fixed texts of lang-es.tex (\renewcommand).
\renewcommand{\tComputerText}{\textbf{Un MSX1, MSX2 o turbo R} con dos ranuras de cartucho, una para el juego y otra para el V9990. En un turbo R el juego va a la misma velocidad que en un MSX1.}
\renewcommand{\tMemoryText}{\textbf{\ramNeeded{} de RAM}. Un MSX con menos memoria también sirve si el cartucho flash que lleva el juego tiene memoria propia (por ejemplo un MegaFlashROM SCC+ SD): el juego la usa.}
\renewcommand{\tSoundText}{El sonido del propio MSX: música del título, sirena de policía, motor, turbo y efectos. No hace falta ningún cartucho de sonido.}
\renewcommand{\tGameText}{El \textbf{cartucho de \gameTitle{} para V9990}: una \romKind{} (el fichero \texttt{\romFile}, grabado en un cartucho flash o equivalente). Funciona con el cartucho en una ranura normal o en un subslot.}
\renewcommand{\tControlsText}{Los cursores y el teclado, o un joystick de dos botones en el puerto 1.}
\renewcommand{\tSetJoy}{Conecta un joystick en el puerto 1 si quieres usarlo. El botón 2 dispara el turbo.}
\renewcommand{\tSetOn}{Enciende el MSX. Durante unos dos segundos la pantalla del propio MSX muestra \texttt{OCEAN SOFTWARE 1989}, \texttt{DIHALT STUDIO 2026} y \texttt{V9990 found !}; después la imagen aparece en el monitor del V9990. Si pone \texttt{V9990 NOT found!!!} es que falta el V9990 o no está bien insertado, y el juego se para ahí.}
% Los créditos van después de las secciones 5-8 del juego.
\renewcommand{\tCreditsTitle}{9. Cr\'editos y derechos}
\newcommand{\tCreditsExtra}{%
\item El \textbf{Chase H.Q.} original es un juego de recreativa de \textbf{Taito Corporation}, 1988. La versión MSX es de Ocean: programa Jobbee, gráficos Bill, música Jon Dunn.
\item La ilustración de la portada es la de la recreativa de Taito. La música y la sirena siguen las de la versión ZX Spectrum 128K.}
